\chapter[Sermon 99. He Gathers Such Things as He Hath Taught of the Last Judgment, and of the Rewards of the Godly, and of the Torments of the Wicked.]{He Gathers Such Things as He Hath Taught of the Last Judgment, and of the Rewards of the Godly, and of the Torments of the Wicked.}
\label{sermon:099}
\markright{Sermon 99. He Gathers Such Things as He Hath Taught of the Last ...}

And behold, I am coming quickly, and my recompense is with me, to give to every man according to his works. I am Alpha and Omega, the beginning and the end, the first and the last. Blessed are those who keep his commandments, that their power may be in the tree of life, and that they may enter through the gates into the city. For outside will be dogs and sorcerers, and adulterers, and murderers, and idolaters, and whoever loves and practices falsehood [? leasings].

The ninth point of this conclusion concerns the Lord’s coming for judgment, and the rewards prepared for the righteous, and the punishments appointed for the unrepentant and wicked. He gathers here that he treated this matter more thoroughly in chapters 19 and 20, and in other passages of this book. And this point, above all others, he emphasizes and urges most earnestly, for it is of great importance if we both understand it correctly and consider it often in our minds. For we will sin less freely, but will watch more diligently.

\index{John, Saint}And in this conclusion of St. John, the speakers are often changed. For now John speaks himself, and immediately he introduces the Lord speaking. Indeed, at this point he makes the Lord Christ himself speak, saying, “Behold, I come quickly.” For the word pronounced from Christ’s mouth has more authority and carries more weight with all than what the Apostle speaks. And in saying that he will come shortly, he intends to stir up all people to watch, repent, and pray. For in the Gospel he said, “Watch, for you know neither the day nor the hour.” Your Lord will come at an hour when you think least. He fears the slothful and unclean people who comfort themselves, believing that the Lord will not come at all, and if he does come, that it will be long first, and perhaps never. Against them, pleading, he says how he will come quickly. Against the same, Malachi reasoned in chapters 3 and 4. And St. Peter reasoned in chapters 2 and 3. Moreover, in affirming that he will come shortly, he comforts the godly who are tempted and tossed diversely in this world. For the godly sometimes cry out that the Lord delays his coming for a long time, that he is too benign to his enemies. Wherefore he says that he will come soon enough, that is to say, in due time, that he may both deliver his servants and destroy and root out his enemies and opponents.

For it follows, what manner of person he is, how, and to what end he will come: he will come gloriously with great majesty and power to deliver and save the faithful, and condemn the ungodly, for he says, “My reward is with me.” These words seem to be taken from the fortieth chapter of Isaiah. And they signify that God is abundantly furnished with all the instruments with which a rewarder and revenger must be furnished. Therefore he says, “The reward which I shall give to every one, after his doings, I have presently with me, and that ready and plentiful.” For our king and judge lacks not power and treasure; as often the kings of this world either cannot pay their soldiers’ wages as they ought, or do not have it ready, and delay the payment a long time. But this our Captain says, “My reward is with me.” And immediately explaining himself, he says that he will reward every one according as his doing shall be. For so the Apostle also in 2 Corinthians 5 says, how we must all appear before the judgment seat of Christ, that every one may receive such things as are done by the body, according as he has done, whether it be good or evil. For in the sixteenth chapter of the gospel of Saint Matthew, the Lord said likewise that there would come a time when the Son of Man should come in the glory of his Father, with his angels, and then he shall render to every one after his doings. The same is taught by the Apostle in Romans 2.

And to ensure that no one should doubt that our judge can accomplish in deed what he said he would do, namely to render to every man according to his works, he adds and says, “I am Alpha and Omega; the beginning and the end,” and so forth. By these words, he signifies that he is truly God, eternal, and almighty. The sentence is taken from the 43rd and 45th chapters of Isaiah and has been explained previously. These things teach us that Jesus Christ is truly God, and therefore the rewarder of all, most bountiful and most righteous.

Consequently again, expressly, more plainly, and by a division, John declares what and to whom the Lord will give. First, indeed, he treats of the reward prepared for the good, then of the punishment appointed for the evil by the just judgment of God. And the reward is paid, or given rather, as Paul says, to those who keep his commandments, namely Christ’s. For not those who read or hear the commandments of God, or boast and preach them, are blessed, but those who keep and perform them in deed. For so our Lord and Savior Christ taught us in the Gospel according to Matthew, chapter 7, and Luke, chapter 11. And his commandments are those that are expounded in the ten precepts, or in the Gospel restrained to the love of God and our neighbor, or which are named by John the Apostle faith and love. It behooves us therefore to be religious, in case we look to receive a reward from God. And what is the reward that is given by the judge to the godly worshippers of God? That is taken in three ways. First, they are called happy and blessed. Secondly, they shall have power over the tree of life, that is to say, the fruits of the tree of life shall be in their power: that is to wit, they shall live an eternal life with Christ, as before is declared. For he alludes to the former things. Lastly, they shall enter in, he says, by the gates into the city (to wit, before also described) into the everlasting country.

After this, he addresses the punishments appointed for the wicked, and in one word encompasses them all, when he says, “without.” For by this single word, he excludes the wicked from the heavenly realm and confines them to hell, where torments are unspeakable, endless, and innumerable. And Saint John here follows the Lord in the Gospel, saying: “I say to you, that many shall come from the East and from the West, and shall rest with Abraham, Isaac, and Jacob in the kingdom of heaven; but the children of the kingdom shall be cast out into the outer darkness, where there will be weeping and gnashing of teeth.” Similarly, in the parable of the ten virgins, the gate is said to be shut, and the foolish virgins shut out from celestial joys. He also commands the unprofitable servant to be cast out into the outer darkness. Likewise, in Luke 13, the Lord says how the unbelievers shall be expelled.

\index{Turks}\index{Repentance}And who, I pray, shall be cast out in that final judgment? Dogs, and those listed in the register of the condemned. The term “dogs” is not always used in Scripture to denote the wicked, but generally it does. Abner, the prince of King Saul’s wars, says to Ishbosheth, “Am I the head of a dog, to defend the house of Saul against Judah?” This signifies that he had incurred the displeasure of the tribe of Judah, because he had retained ten tribes still in their duty and under the dominion of the house of King Saul. Elsewhere, as in Matthew 15, Gentiles, or pagans, or those estranged from the people of God, seem to be called dogs. Some today call the Turks “Turkish dogs,” meaning Turkish unbelievers. Furthermore, the prophet Isaiah calls the false prophets dogs: shameless, ravenous, insatiable, unable to bark and defend the Lord’s flock, or else unwilling and sleepy. With the same meaning, the Apostle says to the Philippians: “Beware of dogs, beware of evil workers, etc.” Moreover, in Scripture, angry, fierce, cruel men, those who despise godly things, who bark at the truth, slanderers, persecutors, and blasphemers are called dogs. For in Psalm 22, David, a figure of Christ the Lord, cries: “Dogs have surrounded me; the counsel of the wicked has encompassed me.” Whom he now calls dogs, he immediately names wicked. And when Shimei cursed David, Abishai the son of Zeruiah says: “Why does this dog, who is to die, curse my lord the king?” And the Lord in the Gospel forbids giving what is holy to dogs, or pearls to swine. Finally, these filthy men, the unclean, are called dogs; they are without repentance, wallowing themselves in the filth of sin and wickedness.

For St. Peter calls such people dogs returning to their vomit. And the Lord forbids that anyone bring the price of a prostitute or a dog into the Temple. Even then, the Jewish priests rejected the price of blood offered by Judas. Therefore, under the name of dogs, we understand pagans or unbelievers, false prophets or deceivers, cruel men, blasphemers, persecutors of the truth, cursed speakers, contemners of the truth, unclean and filthy, and so on.

\index{Primasius}And as for the matters that have been explained before, namely in chapter 9 and at the beginning and end of chapter 21. He adds a lie here, he who loves and promotes. For many do not openly create them, but they love, favor, and advance them. Many both love and create them. They love chiefly a falsehood, which supports false doctrine and delights therein. But concerning this, Primasius, Bishop of Utica, speaks purposefully: he says that to all these things, one must give not diligence of explanation, but carefulness of avoiding the evils. The Lord Jesus save us from all evil. Amen.
